\documentclass{article}
\usepackage[UTF8]{ctex}
\usepackage{amsmath,amssymb,amsfonts}
\usepackage{enumitem}
\usepackage{geometry}
\geometry{a4paper, margin=1in}
% === 错题本排版优化（不参与题目标识） ===
\setlength{\parindent}{0pt}
\setlength{\parskip}{0.35em plus 0.10em minus 0.08em}
\setlength{\emergencystretch}{3em}
\setlist{itemsep=0.25em, topsep=0.40em, parsep=0pt}
\allowdisplaybreaks[2]
% === 错题本排版优化结束 ===
\begin{document}
\section*{河海大学2025年数学分析真题}
\section*{94．河海大学2025年数学分析试卷}

\subsection{题目1}
1．（5 分）数列 \(\left\{a_{n}\right\}\) 为实数列，如果对所有的 \(n, p \in \mathbb{N}^{+}\)，都有 \(\left|a_{n+p}-a_{n}\right|<\frac{p}{n}\) ，那么 \(\left\{a_{n}\right\}\) 为基本数列（Cauchy 列）。

\subsection{题目2}
2．（5 分）如果连续函数 \(f(x)\) 在 \((a, b)\) 上没有零点，那么 \(f(x)\) 在 \((a, b)\) 内不变号．

\subsection{题目3}
3．（5 分）如果函数 \(f(x)\) 在 \([a, b]\) 上可积，那么 \(f(x)\) 在 \([a, b]\) 上至多只有有限个间断点．

\subsection{题目4}
4．（5 分）如果级数 \(\sum_{n=1}^{\infty} a_{n}\) 收敛，实数列 \(\left\{b_{n}\right\}\) 满足 \(\lim _{n \rightarrow \infty} \frac{b_{n}}{a_{n}}=0\) ，那么 \(\sum_{n=1}^{\infty} b_{n}\) 也收敛。

\subsection{题目5}
5．（5 分）写出命题＂当 \(x \rightarrow x_{0}^{-}\)时，函数 \(f(x)\) 的极限为正无穷大＂的否定命题的分析表述．

\subsection{题目6}
6．（5 分）叙述函数 \(f(x)\) 在区间 \(I\) 上一致连续的定义．

\subsection{题目7}
7．（5 分）叙述函数 \(f(x)\) 在 \([a, b]\) 上 Riemann 可积的定义．

\subsection{题目8}
8．（5 分）叙述判定函数项级数 \(\sum_{n=1}^{\infty} a_{n}(x) b_{n}(x)\) 在 \(D\) 上一致收敛的 Dirichlet 判别法．

\subsection{题目9}
9．（10 分）求极限 \(\lim _{x \rightarrow 0} \frac{\frac{x^{2}}{2}-\sqrt{1+x^{2}}+1}{\left(\cos x-e^{x^{2}}\right) \tan ^{2} x}\) ．

\subsection{题目10}
10．（10 分）求不定积分 \(\int \frac{x^{2} \arcsin x}{\sqrt{1-x^{2}}} \mathrm{~d} x\) ．

\subsection{题目11}
11．（10 分）设 \(p>0\) ，讨论级数

\[
\sum_{n=1}^{\infty} \ln \left(1+\frac{(-1)^{n-1}}{n^{p}}\right)
\]

是否收敛？如果收敛的话，是条件收敛还是绝对收敛？

\subsection{题目12}
12．（10 分）函数 \(z=z(x, y)\) 具有二阶连续偏导数，且满足方程

\[
\frac{\partial^{2} z}{\partial x^{2}}-2 \frac{\partial^{2} z}{\partial x \partial y}+\frac{\partial^{2} z}{\partial y^{2}}=0
\]

现做自变量变换 \(u=x-y, v=x+y\) 和因变量变换 \(w=x y-z\) ，求 \(w\) 作为 \(u\) 和 \(v\) 的函数所满足的方程．

\subsection{题目13}
13．（10 分）设 \(h>0\) ，求

\[
\iint_{\sum} x^{3} \mathrm{~d} y \mathrm{~d} z+y^{3} \mathrm{~d} z \mathrm{~d} x+z^{3} \mathrm{~d} x \mathrm{~d} y
\]

其中 \(\sum\) 是雉面 \(z=\sqrt{x^{2}+y^{2}}\) 在 \(0 \leq z \leq h\) 的部分，方向取上侧。

\subsection{题目14}
14．（15 分）函数 \(f(x)\) 在［0，1］上二阶连续可导，且 \(|f(x)| \leq \min \{1-\cos x, \sin \pi x\}, \forall x \in[0,1]\) ．证明：

\begin{itemize}
  \item[（1）] 存在 \(\xi \in(0,1)\) ，使得 \(f^{\prime}(\xi)=0\) ．
  \item[（2）] 存在 \(\eta \in(0,1)\) ，使得 \(f^{\prime \prime}(\eta)=0\) ．
\end{itemize}

\subsection{题目15}
15．（15 分）设 \(f(x, y)=\frac{1}{(1+x y)(1+x)},(x, y) \in[0,+\infty) \times[0,+\infty)\) ．

\begin{itemize}
  \item[（1）] 证明 \(\int_{0}^{+\infty} f(x, y) \mathrm{d} x\) 关于 \(y \in(0,+\infty)\) 上是内闭一致收敛的．
  \item[（2）] 设 \(g(y)=\int_{0}^{+\infty} \frac{\ln (1+x y)}{x(1+x)} \mathrm{d} x, y \in(0,+\infty)\) ，求 \(g^{\prime}(y)\) 及 \(g(1)\) ．
\end{itemize}

\subsection{题目16}
16．（15 分）解答如下问题：

\begin{itemize}
  \item[（1）] 求幂级数 \(\sum_{n=1}^{\infty} \frac{x^{n}}{n}\) 的收敛域以及和函数。
  \item[（2）] \(f(x)=\sum_{n=1}^{\infty} \frac{x^{n}}{n^{2}}, x \in[-1,1]\) ．证明：对任意的 \(x \in(0,1)\) ，都有
\[
f(x)+f(1-x)+\ln x \cdot \ln (1-x)=\frac{\pi^{2}}{6}
\]
\end{itemize}

\subsection{题目17}
17．（15 分）设数列 \(\left\{a_{n}\right\}\) 满足 \(\lim _{n \rightarrow \infty} a_{n}=a\) ，数列 \(\left\{p_{n}\right\}\) 满足 \(p_{k}>0, k=1,2, \cdots\) ，且

\[
\lim _{n \rightarrow \infty} \frac{p_{n}}{p_{1}+p_{2}+\cdots+p_{n}}=0
\]

证明：

\[
\lim _{n \rightarrow \infty} \frac{p_{1} a_{n}+p_{2} a_{n-1}+\cdots+p_{n} a_{1}}{p_{1}+p_{2}+\cdots+p_{n}}=a .
\]
\end{document}
